\section{Marco Normativo}

\subsection{Ventilación y Renovaciones de Aire}

ANSI/ASHRAE 62.1 como referencia de calidad de aire interior y ANSI/ASHRAE 170 (ventilación de instalaciones de salud) como referencia de tasas de renovación y niveles de presurización; las 12 renovaciones/h adoptadas superan holgadamente los mínimos de ambas para el uso previsto (dato típico de aplicación).

\subsection{Filtración}

ANSI/ASHRAE 52.2-2017 (clasificación MERV, base de la especificación) e ISO 16890 (equivalencia ePM1 exigible en informes de ensayo de proveedores internacionales); EN 779:2012 como referencia heredada (F7/F8).

\subsection{Ventiladores y Dampers}

AMCA 210/211 (ensayo y certificación de desempeño), AMCA 500-D (caída de presión y fuga de dampers), AMCA 99 (construcción Spark A), ASTM C582 y ASTM D4167 (laminados PRFV), ASHRAE 70 (medición de rejillas y difusores).

\subsection{Instalación Eléctrica}

RETIE (Reglamento Técnico de Instalaciones Eléctricas, Colombia) y NTC 2050 (Código Eléctrico Colombiano) para acometida, protecciones, puesta a tierra y certificación de la instalación del motor y el tablero de control; NEMA MG-1 / IEC 60034-1 para derateo del motor por altitud; IEEE 841 como referencia de ejecución severe duty.

\subsection{Seguridad Ocupacional}

OSHA 29 CFR 1910.1450 (exposición ocupacional a químicos peligrosos en laboratorios) como referencia de buena práctica de ventilación de laboratorios, complementada por la Resolución 0312 de 2019 (estándares mínimos de seguridad y salud en el trabajo, Colombia).

\subsection{Protección Anticorrosiva}

ISO 12944 (pinturas y recubrimientos; el ambiente corresponde a categoría de corrosividad alta, C5, por la combinación cloro/humedad, dato típico de clasificación) y prácticas NACE/AMPP para selección de materiales en servicio clorado.
